\cotefile{158}
\scanpage{20}
\mspage{1}

\uline{Proposition 1} Soit $X\in\mathrm{Ob}\,\mathrm{Cat}$ \deleted{\unsure{et $0$-co}} une catégorie pseudo-cofiltrante ($\Longleftrightarrow X^\wedge$ est localement irréductible), ayant un plus petit objet $e$\inserted{(donc $X$ est \unsure{$0$-conn.|$0$-connexe}).}

Soit $X_0$ la sous-catégorie pleine réduite à $e$, $M=\mathrm{End}_X(e)=\mathrm{End}_{X_0}(e)$ le monoïde des endomorphismes de $e$, $G=\pi_1(M)=\pi_1(X_0,e)$ le groupe enveloppant de $M$. Alors

a) L'hom canonique $\pi_1(X_0,e)\longrightarrow\pi_1(X,e)$ est surjectif, \deleted{et \unsure{son|pour}}

b) Pour toute double-flèche de source $e$
\diagram{\begin{tikzcd}[ampersand replacement=\&]e\arrow[r,shift left=.55ex,"u"]\arrow[r,shift right=.55ex,"v"']\&x\end{tikzcd}}
dans $X$, il existe $u',v'\in M$ tels que
\[
uu'=vv',
\]
et l'élément
\[
\unsure{[v]^{-}[u]|[v]^{-1}[u]|[v][u]}\ \text{de}\ \pi_1(X,e)
\]
(image canonique du chemin \diagram{\begin{tikzcd}[ampersand replacement=\&]e\arrow[r,"u"]\&x\&e\arrow[l,"v"']\end{tikzcd}})
est égale à l'image dans $\pi_1(X,e)$ de l'élément $[v,u]\overset{\mathrm{def}}{=}v'u'^{-1}\in G=\pi_1(X_0,e)$.

De plus, \inserted{si $e$ est aussi objet maximal de $X$,} cet élément \inserted{$[v,u]$} de $G$ ne dépend pas du choix du couple $(u',v')$ \deleted{\unsure{d'}}d'éléments de $M$ satisfaisant $uu'=vv'$. Plus précisément, si $\alpha:x\longrightarrow e$, alors $\deleted{\unsure{w u=v=}}u_1=\alpha u$ et $v_1=\alpha v$ sont des éléments $\deleted{[v,u]=}$ de $M$, et on a $u_1u'=v_1\unsure{u'|v'}$ d'où
\[
[v,u]=v_1^{-1}u_1^{\deleted{\unsure{-1}}}.
\]

c) Pour toute double-flèche \diagram{\begin{tikzcd}[ampersand replacement=\&]e\arrow[r,shift left=.5ex,"u"]\arrow[r,shift right=.5ex,"v"']\&x\end{tikzcd}}, dans $X$,

% archive sheet number: 19
